\section{Does \pizf See the Displacement, and Does It Use It?}
\label{sec:probing}

\paragraph{Model.}
\pizf~\cite{pi2025pi05} couples a PaliGemma VLM~\cite{beyer2024paligemma} (18 layers, width 2048) with an action expert
(18 layers, width 1024) that produces a 50-step action chunk by flow matching~\cite{lipman2023flow}.
At every layer the action tokens attend to the VLM's prefix: 512 image tokens from an agent-view and a wrist camera, the
instruction, and the proprioceptive state, which the LIBERO checkpoint writes into the prompt as about 120 digit tokens.
The time step of the flow conditions the expert through adaptive RMS normalization.
We use the public LIBERO checkpoint in its LeRobot port; on LIBERO-PRO it matches the numbers reported for the official
model (swap: Object 18.7 \vs 18.2, Long 8.5 \vs 9.4~\cite{chen2026ect}).

\paragraph{Decision-time interventions.}
We take states $s_t$ from successful rollouts of the policy and build a counterfactual state $s'_t$ that differs in one
respect, by editing the simulator's kinematic state and re-rendering both cameras: (i) before the grasp, the target is
moved by 3, 6, 12 or 20\,cm to a free spot, or (ii) it swaps places with another object; (iii) while the object is
carried, the container is moved by the same amounts; (iv) before the grasp, the instruction names another object of the
scene; (v) only the proprioceptive input is offset by 5\,cm.
Nothing else changes, the robot included.
States are binned by phase: \emph{early} (hand farther than 15\,cm from the target), \emph{mid} (5--15\,cm), \emph{near}
(closer than 5\,cm) and \emph{carrying}.

\paragraph{Use.}
With the flow noise held fixed, we roll the first 25 steps of each predicted chunk through a kinematic model of the
controller to obtain the hand's intended displacement $m(s)$, and define
\begin{equation}
\mathrm{use} = \frac{\langle m(s'_t) - m(s_t),\, \mathbf{d} \rangle}{\lVert \mathbf{d} \rVert^2},
\end{equation}
where $\mathbf{d}$ is the displacement of what moved (for (iv), from the target to the newly named object).
A value of 1 means the motion follows the change fully, 0 that it ignores it.

\paragraph{Encoding.}
For every layer we pool the hidden states of four streams (image tokens, instruction tokens, the last prefix token and the
action tokens) and regress the \emph{within-state} difference $h(s'_t) - h(s_t)$ on $\mathbf{d}$ with ridge regression,
reporting cross-validated $R^2$ with folds grouped by episode.
Differencing the same state removes everything that varies across states (where the hand is, the time in the episode),
so the probe only asks whether the stream registers where the object was moved to.

\begin{table}[t]
\centering
\small
\setlength{\tabcolsep}{4pt}
\begin{tabular}{lccccc}
\toprule
Phase & 3\,cm & 20\,cm & swap & rename & proprio \\
\midrule
early ($>$15\,cm) & 0.18 & 0.22 & 0.22 & 0.25 & 0.00 \\
mid (5--15\,cm) & 0.61 & 0.31 & \textbf{0.03} & \textbf{0.02} & 0.00 \\
near ($<$5\,cm) & 0.51 & 0.29 & 0.01 & $-$0.02 & 0.00 \\
carrying & \textbf{0.09} & \textbf{0.06} & -- & -- & 0.00 \\
\bottomrule
\end{tabular}
\caption{\textbf{Use of a displacement by \pizf}, LIBERO-Object (median over states). Columns: target moved by 3 or
20\,cm (container while carrying), target swapped with another object, target renamed in the instruction, proprioception
offset by 5\,cm. At the decision moments (bold) the action barely follows the change.}
\label{tab:use_base}
\end{table}

\paragraph{Findings.}
\Cref{tab:use_base} reads as follows.
\emph{(1) At the decision moments the policy is nearly blind.}
Early in the approach the action follows a moved target by about 20\%; mid-way it ignores a swap (0.03) or a renamed
target (0.02) entirely, and while carrying it follows a moved container by 6--9\%: placing is a memorized motion.
\emph{(2) Once close, it is a local tracker.}
Small displacements near the hand are followed (0.51--0.61) more than large ones, but the choice of object is locked.
\emph{(3) Proprioception is not the shortcut.}
Offsetting the proprioceptive input has no effect in any phase, and replacing it by a constant at test time (Object swap
14.0, nominal 100) or during training (Spatial swap 40.5 \vs 40.0) does not change the failure.
The same pattern holds on LIBERO-Spatial (early 0.21--0.30, swap 0.08--0.14, carrying 0.17--0.30) and LIBERO-Goal.

At the same time, the displacement is \emph{encoded}.
It is decodable from the image tokens at every layer ($R^2$ 0.7--0.8 for a target moved by 12\,cm or more, about 0.9 for a
moved container) and is carried into the action tokens over the expert's layers, from about 0.05 in its first layers to
0.58 (target) and 0.74--0.82 (container) in its last ones.
The attention of the action tokens is also nearly identical between blind and tracking phases (instruction 0.6--0.8 at a
few layers, wrist image 0.2--0.3, agent image 0.1--0.3 in early layers), so the failure is not that the action tokens do
not look at the image.
\pizf knows where the target and the container are; at the moments that matter, its action does not depend on it.
